\documentclass{article}
\usepackage[landscape, margin=0.28cm]{geometry}
\usepackage{multicol}
\usepackage{amssymb}
\usepackage{amsmath}
\usepackage{enumitem}
\usepackage[compact]{titlesec}
\usepackage{lmodern}
\usepackage{mathpazo}
\renewcommand{\familydefault}{\sfdefault}
\usepackage{bm}
\usepackage{ctex}
\usepackage{xcolor}
\pagenumbering{gobble}
\definecolor{mycolor}{RGB}{200, 0, 0} 

\titleformat{\section}{\normalfont\fontsize{9.2pt}{10pt}\bfseries\color{mycolor}}{\thesection}{0.18em}{}
\titleformat{\subsection}{\normalfont\fontsize{8.2pt}{9pt}\bfseries\color{mycolor}}{\thesubsection}{0.18em}{}
\titleformat{\subsubsection}{\normalfont\fontsize{7.4pt}{8pt}\bfseries\color{mycolor}}{\thesubsubsection}{0.18em}{}
\titlespacing*{\section}{0pt}{2pt}{1pt}
\titlespacing*{\subsection}{0pt}{1.5pt}{0.5pt}
\titlespacing*{\subsubsection}{0pt}{1pt}{0pt}
\setlength{\columnsep}{0.32cm}
\setlength{\columnseprule}{0.18pt}
\setlength{\parindent}{0pt}
\setlength{\parskip}{0pt}
\usepackage{setspace}
\setstretch{0.9}

\setlist[itemize]{nosep, leftmargin=*, label=\textbullet}
\setlist[enumerate]{nosep, leftmargin=*}
\allowdisplaybreaks
\setlength{\abovedisplayskip}{0.5pt}
\setlength{\belowdisplayskip}{0.5pt}

\newcommand{\key}[1]{\textbf{#1}}
\newcommand{\lowterm}[1]{\textbf{\underline{#1}}}
\newcommand{\term}[1]{\textit{#1}}
\newcommand{\R}{\mathbb{R}}
\newcommand{\E}{\mathbb{E}}
\newcommand{\N}{\mathcal{N}}
\newcommand{\KL}{\mathrm{KL}}
\newcommand{\1}{\mathbf 1}
\newcommand{\dd}{\mathrm{d}}
\DeclareMathOperator*{\argmax}{arg\,max}
\DeclareMathOperator*{\argmin}{arg\,min}
\DeclareMathOperator{\softmax}{softmax}
\DeclareMathOperator{\sigmoid}{sigmoid}
\DeclareMathOperator{\tr}{tr}

\begin{document}
\tiny
\begin{multicols*}{4}

\section{L2--3 Supervised Learning}
\subsection{Classification and Loss}
\key{ERM}: population risk $\theta^\star=\arg\min_\theta\E_{x\sim P}\operatorname{err}(f(x;\theta),y(x))$; empirical risk $\hat\theta=\arg\min_\theta\frac1N\sum_i\operatorname{err}(f(x_i;\theta),y_i)$. \key{Binary logistic}: $p(y=1|x)=\sigma(w^\top x+b)$,
\[
\ell=-y\log p-(1-y)\log(1-p),\quad \frac{\partial\ell}{\partial z}=p-y.
\]
\key{Multi-class}: logits $z\in\R^K$, $p_k=e^{z_k}/\sum_j e^{z_j}$,
\[
\ell(z,y)=-\log p_y=-z_y+\log\sum_j e^{z_j},\quad
\frac{\partial\ell}{\partial z_k}=p_k-\1[k=y].
\]
\lowterm{One-hot CE}: $-\sum_k y_k\log p_k$.

\subsection{MLP and Backprop}
\key{Layer}: $h^{(0)}=x$, $z^{(\ell)}=W^{(\ell)}h^{(\ell-1)}+b^{(\ell)}$, $h^{(\ell)}=\phi(z^{(\ell)})$. Batch code convention often uses $X\in\R^{B\times n}$, $Z=XW+b$, $W\in\R^{n\times m}$.
\key{Backprop}: define $\delta^{(\ell)}=\partial L/\partial z^{(\ell)}$:
\(
\partial_{W^{(\ell)}}L=\delta^{(\ell)}(h^{(\ell-1)})^\top,\quad
\partial_{b^{(\ell)}}L=\delta^{(\ell)},\quad
\)
\(
\delta^{(\ell-1)}=((W^{(\ell)})^\top\delta^{(\ell)})\odot\phi'(z^{(\ell-1)}).
\)

\key{Activation}: sigmoid/tanh saturate; ReLU $\max(0,z)$ keeps positive-side gradient; ELU/smooth variants improve negative side. \lowterm{Subgradient}: ReLU at $0$ uses any value in $[0,1]$ by convention.

\subsection{CNN}
\key{Motivation}: locality + weight sharing. Convolution:
\[
y_{i,j,o}=\sum_{u,v,c}K_{u,v,c,o}x_{i+u,j+v,c}+b_o.
\]
For input $H\times W\times C_{\rm in}$, kernel $K_h\times K_w$, padding $P$, stride $S$:
\(
H'=\lfloor(H+2P-K_h)/S\rfloor+1,\quad
W'=\lfloor(W+2P-K_w)/S\rfloor+1.
\)
Params $K_hK_wC_{\rm in}C_{\rm out}+C_{\rm out}$. \lowterm{Scanning MLP}: same detector slides across space. \lowterm{Pooling}: max/avg reduces spatial size and jitter. \lowterm{Padding}: keeps boundary info; same padding for odd kernel uses $P=(K-1)/2$.

\subsection{Optimization}
\lowterm{GD}: $x_{k+1}=x_k-\eta_k\nabla f(x_k)$. \lowterm{$L$-smooth}: $\|\nabla f(x)-\nabla f(y)\|\le L\|x-y\|$; if $C^2$, $\nabla^2f(x)\preceq LI$. It implies
\(
f(y)\le f(x)+\nabla f(x)^\top(y-x)+\frac L2\|y-x\|^2,
\)
so $0<\eta<2/L$ decreases the upper bound. 
\lowterm{Strong convexity}: $f(y)\ge f(x)+\nabla f(x)^\top(y-x)+\frac\mu2\|y-x\|^2$; with $L$-smooth and $C^2$, $\mu I\preceq\nabla^2f(x)\preceq LI$. GD has linear convergence $k=O(\kappa\log(1/\epsilon))$, $\kappa=L/\mu$.

\lowterm{Newton}: from quadratic Taylor,
\(
x_{k+1}=x_k-\eta[\nabla^2f(x_k)]^{-1}\nabla f(x_k).
\)

$g_k=$ gradient on $x_k$. 
\key{Momentum}: $\Delta X^k=\beta\Delta X^{k-1}-\eta\nabla f(X^k)$, $X^{k+1}=X^k+\Delta X^k$.

\key{Nesterov}: compute gradient at lookahead $x_k+\beta(x_k-x_{k-1})$.

\key{SGD / mini-batch}: use one or a small batch of samples to estimate the full gradient; mini-batch reduces variance while keeping updates cheap. \key{Adaptive LR}: per-coordinate step sizes help convergence. \key{AdaGrad}: $s_k=s_{k-1}+g_k^2$, $x_{k+1}=x_k-\eta g_k/(\sqrt{s_k}+\epsilon)$; rarely updated dimensions have smaller $s_{k,i}$ and larger effective learning rate, but learning can decay too fast. 

\key{RMSProp}: $s_k=\rho s_{k-1}+(1-\rho)g_k^2$. 

\key{Adam}:
\(
m_k=\beta_1m_{k-1}+(1-\beta_1)g_k,\quad v_k=\beta_2v_{k-1}+(1-\beta_2)g_k^2,
\)

\(
x_{k+1}=x_k-\eta\hat m_k/(\sqrt{\hat v_k}+\epsilon),\quad
\hat m_k=m_k/(1-\beta_1^k).
\)

\subsection{Regularization and Architectures}
\lowterm{Weight decay}: L2 regularization stabilizes weights, $L_{\rm reg}(w)=L_{\rm loss}(w)+\alpha\|w\|_2^2$. \lowterm{Early stopping}: stop when training improves but held-out validation worsens. \lowterm{Dropout}: randomly zero units, prevents overfitting to particular neurons/weights, approximates ensemble and reduces co-adaptation; can be viewed as implicit L2. Two conventions: scale activations by keep-prob during training and leave inference unchanged, or leave training unchanged and multiply by keep-prob at inference. \lowterm{Gradient clipping}: $g\leftarrow g\min(1,\tau/\|g\|)$.

\lowterm{Normalization}: \key{GroupNorm}: batch-independent. \key{InstanceNorm}: batch-independent, suitable for generation tasks. 
\key{BatchNorm}: $\mu_B=\frac1B\sum_i z_i$, $\sigma_B^2=\frac1B\sum_i(z_i-\mu_B)^2$,
\(
\hat z_i=\gamma\frac{z_i-\mu_B}{\sqrt{\sigma_B^2+\epsilon}}+\beta.
\)
Inference uses running averages because there is no mini-batch training distribution: $\mu_{\rm run}\approx\frac1{N_{\rm batch}}\sum_{\rm batch}\mu_B$, $\sigma_{\rm run}^2\approx\frac1{N_{\rm batch}}\sum_{\rm batch}\frac{B}{B-1}\sigma_B^2$.
Backprop through a layer uses Jacobian transpose: $\nabla_xL=J^\top\nabla_yL$.
\key{PCA color jitter}: $[I_R,I_G,I_B]\leftarrow[I_R,I_G,I_B]+\alpha[P_R,P_G,P_B][\lambda_R,\lambda_G,\lambda_B]^\top$, $\alpha\sim\N(0,0.1)$; perturb RGB along principal color directions.

\key{LayerNorm}: normalize hidden dimension inside one sample; better for sequence models.  Batch-independent, Particularly suitable for RNN. $h^t=f\left(\frac{g}{\sigma^t}\odot(a^t-\mu^t)+b\right), \mu^t=\frac{1}{H}\sum_i a_i^t, \sigma^t=\sqrt{\frac{1}{H}\sum_i (a_i^t-\mu^t)^2}$.  


\key{ResNet}: $h_{\ell+1}=h_\ell+F(h_\ell)$, residual path eases optimization. \key{DenseNet}: concatenate previous features. \key{FCN}: replace fully connected layer by convolution for dense prediction.

\section{L4 Energy-Based Models}
\subsection{Energy View}
\key{Generative modeling}: learn $P(X)$ or $P(X,y)=P(y)P(X|y)$ instead of $P(y|X)$.  Generative models can be used to both classify and generate images

\lowterm{EBM}: unnormalized score via energy
\[
p_\theta(x)=\frac{e^{-E_\theta(x)}}{Z_\theta},\quad Z_\theta=\int e^{-E_\theta(x)}dx.
\]
Low energy means high probability; training and inference both need sampling or partition-function gradients. $\log p_\theta(x)=-E_\theta(x)-\log Z_\theta$, $\nabla_\theta\log p_\theta(x)=-\nabla_\theta E_\theta(x)-\frac{1}{Z_\theta}\nabla_\theta \left(\sum_{x'} e^{-E_\theta(x')}\right)=-\nabla_\theta E_\theta(x)+\E_{x'\sim p_\theta}[E_\theta(x')]$. 

We can train an energy-based model without knowing the explicit density function
(or normalizing factor Z). Such as score-matching, etc.

\subsection{Hopfield Network}
State $y_i\in\{-1,1\}$, local field $h_i=\sum_{j\ne i}w_{ij}y_j+b_i$. Update flips if $y_ih_i<0$.
\[
E(y)=-\sum_{i<j}w_{ij}y_iy_j-\sum_i b_iy_i=-\frac12y^\top Wy-b^\top y.
\]
Asynchronous flip changes energy by $\Delta E=2y_ih_i$; if $y_ih_i<0$, energy decreases. Finite states imply convergence to a local minimum. \key{Hebbian storage}: $w_{ij}=\frac1{N_p}\sum_p y_i^py_j^p$, $w_{ii}=0$. For $N$ neurons, $K\le N/(4\log N)$ random patterns can be stable. Multiple patterns cause \lowterm{spurious minima}; capacity is limited.

\lowterm{Objective for $W$}: $W=\arg\min_W\left(\sum_{y\in P} E(y)-\sum_{y'\notin P}E(y')\right)$, $E(y)=-\frac12y^\top Wy$. 
\key{Optimization view}: desired patterns should be valleys; short-run evolution can raise undesired valleys by contrastive updates. 

SGD update rule for $W$: $W\gets W+\eta\left(\E_{y\in P}[yy^\top]-\E_{y'}[y'y'^\top]\right)$; in practice sample desired $y$, initialize $y'$ from a desired pattern, run $2$--$4$ timesteps, update with $yy^\top-y'y'^\top$.

\subsection{Boltzmann Machine}
\key{Stochastic Hopfield}: energy $E(y)$ as Hopfield, $p(y)=\frac{1}{Z}e^{-E(y)}\propto e^{-E(y)}$. 

Conditional Gibbs update for $\{-1,1\}$ units: for $h_i=\sum_jw_{ij}y_j+b_i$,
\[
P(y_i=1|y_{\neg i})=\frac{1}{1+e^{-2h_i}}
=\sigma(2h_i),
\quad h_i=\sum_jw_{ij}y_j+b_i.
\]
\[
\log\frac{P(y_i=1|y_{\neg i})}{1-P(y_i=1|y_{\neg i})}=2h_i.
\]

Stochastic evolution: initialize $y^{(0)}$, sample coordinates $y_i(t+1)\sim{\rm Bernoulli}(\sigma(z_i(t)))$, repeat; recover pattern by last-$M$ majority $\hat y_i=\1[\frac1M\sum_{t=L-M+1}^Ly_i(t)>0]$. Temperature $T$: $y_i(T)\sim{\rm Bernoulli}(\sigma(z_i(T)/T))$, $T\leftarrow\alpha T$; annealing lowers $T$ from exploration to exploitation. 

\lowterm{BM Training}: maximize $L(W)=\sum_{y\in P}\log P_\theta(y)$, $\theta$ means  parameters. 
\[
P(y)=\frac{\exp(\frac12y^\top Wy)}{\sum_{y'}\exp(\frac12y'^\top Wy)},\quad
L(W)=\frac1{N_P}\sum_{y\in P}\frac12y^\top Wy-\log\sum_{y'}e^{\frac12y'^\top Wy}.
\]
\(
\nabla_{w_{ij}}\log p_\theta(y)=y_iy_j-\E_{p_\theta}[y_iy_j],
\)
$\nabla_{w_{ij}}L=\frac1{N_P}\sum_{y\in P}y_iy_j-\frac1M\sum_{y'\in S}y_i'y_j'$. 
First term: compute average over $P$. Second term: sample from $p_\theta$, compute average. \key{Sample method}: random initialize $y(0)$, cycle over $i$ (coordinates) and sample from $P(y_i(t)|y_{\lnot i}(t))$, after convergence select the final $M$ states as samples. \key{Updating}: $w_{ij}\gets w_{ij}+\eta \nabla_{w_{ij}}L(W)$.
Learning = \lowterm{positive phase} data correlation minus \lowterm{negative phase} model correlation. 

\lowterm{With hidden neurons}: $y=(v,h)$
\[
\nabla_{w_{ij}} L=\frac{1}{|P|}\sum_{v\in P} \E_{h|v}[y_iy_j]-\E_{y'\sim p_W}[y'_iy'_j].
\]
Use conditioned samples $S_c$ by clamping visible $v$ and Gibbs-sampling hidden $h$, and free samples $S$ by running all variables: $\nabla_{w_{ij}}L=\frac1{N_PK}\sum_{y\in S_c}y_iy_j-\frac1M\sum_{y'\in S}y_i'y_j'$.

Boltzmann for classification: $y=(v,h,c)$, $c$: visible, one-hot vector


\subsection{RBM and Contrastive Divergence}
\lowterm{RBM}: bipartite $v,h$; no visible-visible or hidden-hidden edges:
\(
E(v,h)=-v^\top Wh-b^\top v-c^\top h.
\)
Conditionals factorize:
\(
P(h_j=1|v)=\sigma(c_j+\sum_iw_{ij}v_i),\quad
P(v_i=1|h)=\sigma(b_i+\sum_jw_{ij}h_j),\quad \sigma(z)=\frac{1}{1+e^{-z}}
\). Sampling: $v^{(0)}\to h^{(0)}\to v^{(1)}\to h^{(1)}\to\cdots$

\lowterm{Maximum likelihood estimate}: $\nabla_{w_{ij}}L=\frac{1}{N_PK}\sum_{v\in P}v_i^{0}h_j^{0}-\frac{1}{M}\sum v_i^{\infty}h_j^{\infty}$. CD-1 approximation:
\[
\nabla_{w_{ij}}L\approx \frac1{N_P}\sum_{v\in P}(v_i^0h_j^0-v_i^1h_j^1).
\]

\key{CD-$k$}: start at data $v^{(0)}$, run $k$ Gibbs steps, update $W\propto v^{(0)}h^{(0)\top}-v^{(k)}h^{(k)\top}$.

\subsection{Sampling}
\key{Inverse CDF}: sample $u\sim U[0,1]$, choose category by cumulative mass. \key{Box-Muller}: Gaussian from uniform:
\[
z_1=\sqrt{-2\log u_1}\cos(2\pi u_2),\quad z_2=\sqrt{-2\log u_1}\sin(2\pi u_2).
\]
\lowterm{Importance sampling}: $\E_p[f(x)]=\E_q[f(x)p(x)/q(x)]$. Good $q$ must cover $p$'s high-mass region.
\lowterm{MCMC}: \key{Detailed balance} $\pi(x)T(x\to x')=\pi(x')T(x'\to x)$ implies stationarity of $\pi$. \key{Ergodic}: can visit any desired state with positive probability and has a unique stationary distribution; a strong minorization form is $\min_{z,z':\pi(z')>0}T(z\to z')/\pi(z')=\delta>0$. \key{Metropolis Hastings Algo.}: Construct Markov chain with stationary $\pi$ (desired distribution). Current $x$, draw $x'\sim q(x'|x)$ (called proposal), transition to $x'$ with prob. 
\(
\alpha=\min\left(1,\frac{\pi(x')q(x|x')}{\pi(x)q(x'|x)}\right),
\)
otherwise stays at $x$. 
Choice of $q$: random proposal $q(x'|x)=$noise (e.g. Gaussian, then it relies on $|x-x'|$), then $q$ terms in $\alpha$ vanishes. 

\key{For EBM, partition $Z$ cancels. }

\lowterm{Gibbs}: sample a single coordinate per step from $q(s_i\to s_i')=p(s_i'|s_{j\ne i})$; it is MH with acceptance probability $1$. 

\key{SGLD}: $x_{t+1}=x_t-\frac{\eta}{2}\nabla_xE_\theta(x_t)+\sqrt{\eta}\xi_t$.

\section{L5 Generative Adversarial Networks}
\subsection{Implicit Generative Model}
Instead of density $p_\theta(x)$, learn sampler $x=G_\theta(z)$, $z\sim p_z$. Likelihood-free learning. Need differentiable distance between generated and real distributions. GAN learns this distance with a discriminator.

\lowterm{GAN}: generator $G_\theta(z)$, $z\sim p(z)=N(0,I)$; discriminator $D_\phi(x)$, classify the data is from $p_{data}$ or $G_\theta$. Objective: $L=\min_\theta\max_\phi V(D_\phi,G_\theta)$, 
\[
V(D_\phi,G_\theta)=\E_{x\sim p_{\rm data}}\log D_\phi(x)+\E_{z\sim p_z}\log(1-D_\phi(G_\theta(z))).
\]
Dataset: $\{(x,1):x\sim p_{data}\}\cup\{(\hat{x},0):\hat{x}\sim G_\theta(z)\}$.

Train $D$: $L(\phi)=\E_{x\sim p_{\rm data}}\log D_\phi(x)+\E_{z\sim p_z}\log(1-D_\phi(G(z)))$. 

Train $G$: $L(\theta)=\E_{z\sim p(z)}[\log D_\phi(G_\theta(z))]$.

For fixed $G$, the optimal $D$ is 
\(
D^\star(x)=\frac{p_{\rm data}(x)}{p_{\rm data}(x)+p_G(x)}.
\) 
Substitute:
\(
V(D^\star,G)=-\log4+2\,{\rm JSD}(p_{\rm data}\|p_g), 
\)
\lowterm{JSD}: ${\rm JSD}(p\|q)={\rm KL}(p\|\frac{p+q}{2})+{\rm KL}(q\|\frac{p+q}{2})$, $\geq 0$, $\sqrt{\rm JSD}$ satisfies triangular inequality. 

$2{\rm KL}(p\|q)\geq \left(\int |p(x)-q(x)|\dd x\right)^2$. 

Non-saturating generator loss $-\E_z\log D(G(z))$ gives stronger gradients.

\subsection{Evaluation and Failure Modes}
\lowterm{Mode collapse}: generator covers few modes. For multimodal $p$, reverse KL $\KL(q\|p)=\E_q\log(q/p)$ and JSD often prefer one high-quality mode; forward KL $\KL(p\|q)=\E_p\log(p/q)$ penalizes missing modes and is mode-covering. JSD: $\frac12\KL(p\|m)+\frac12\KL(q\|m)$, $m=(p+q)/2$. \lowterm{Training instability}: moving target game can oscillate; discriminator too strong/simple $\Rightarrow$ generator gradient vanishes.

\key{Inception Score}:
\[
{\rm IS}=\exp\left(\E_{x\sim G}\KL(f(y|x)\|p_f(y))\right),\quad p_f(y)=\E_{x\sim G}f(y|x).
\]
Since $\E_x\KL(f(\cdot|x)\|p_f)=-\E_xH(f(\cdot|x))+H(p_f)$, high IS needs confident images and diverse classes. IS never uses $p_{\rm data}$, so it cannot detect memorization. \lowterm{FID}: Compute $\mu_p,\Sigma_p$ and $\mu_G,\Sigma_G$ for $p_{data},G(z)$ using Inception-v3 pool3 layer (2048-d); compare Gaussian feature statistics:
\(
\|\mu_p-\mu_G\|_2^2+\tr(\Sigma_p+\Sigma_G-2(\Sigma_p\Sigma_G)^{1/2}).
\). Lower FID score means better generators.

\subsection{GAN Training Techniques}
\key{Feature Matching}: train $G$ to match discriminator feature statistics:
\[
L_G(\theta)=\left\|\E_{\hat x\sim G}f(\hat x;\phi)-\E_{x\sim p_{\rm data}}f(x;\phi)\right\|_2^2,
\]
where $f(x;\phi)$ is the final activation layer of $D$. \key{Other tricks}: One-Sided Label Smoothing; Historical Averaging; Minibatch Discrimination; Virtual Batch Normalization.

\subsection{WGAN Family}
JSD can be flat when supports do not overlap. Wasserstein-1:
\(
W(p,q)=\sup_{\|f\|_L\le1}\E_{x\sim p}f(x)-\E_{x\sim q}f(x).
\) 

\lowterm{WGAN}: critic $f_\phi$, no $D$. 
WGAN objective:
\[
\min_G\max_{D:|D|_L\le1}\E_{x\sim p_{\rm data}}[D(x)]-\E_{\hat x\sim p_G}[D(\hat x)].
\]
Equivalently maximize critic $L(\phi)=\E_{x\sim p_{\rm data}}f_\phi(x)-\E_{\hat x\sim p_G}f_\phi(\hat x)$ and train $G$ to reduce the critic gap.


\lowterm{Lipschitz constraint}: weight clipping, spectral norm, or \key{WGAN-GP}: 
$\tilde{x}\gets (1-\epsilon)x+\epsilon\hat{x}$, $\epsilon\sim {\rm Unif}(0,1)$, 
\(
L(\phi)=\E_{x\sim p_{data}}[f_\phi(x)]-\E_{\hat{x}\sim p_G}[f_\phi(\hat x)]+\lambda\E_{\tilde{x}}(\|\nabla_{\tilde x}f_\phi(\tilde x)\|_2-1)^2.
\) 

\key{DCGAN}: conv/deconv, batchnorm, stable architecture. \key{BigGAN}: scale + class conditioning + truncation trick. \key{GigaGAN}: text-to-image GAN scaling. \key{R3GAN}: regularized robust objective.

\subsection{Variants and Applications}
\key{Conditional GAN}: $G(z,c)$, $D(x,c)$. \key{Semi-supervised GAN}: discriminator has $K+1$ classes, $K$ real labels plus fake. \key{BiGAN/ALI}: learn encoder $E(x)$ with discriminator over $(x,z)$ pairs. Idea: $(z,G(z))$ and $(E(x),x)$ from same distribution. gives representation learning. \key{Pix2Pix}: paired image-to-image translation with conditional GAN + reconstruction loss. \key{CycleGAN}: unpaired translation using cycle consistency
\(
G:X\to Y,\quad F:Y\to X,\quad \|F(G(x))-x\|+\|G(F(y))-y\|.
\)
\key{DragGAN}: interactive image manipulation by optimizing latent/code with handle-point constraints.

\section{L6 Flow and VAE}
\subsection{Normalizing Flow}
Invertible differentiable generator $x=G_\theta(z)$,  transform a simple distribution to the data distribution. 
\(
\log p_X(x)=\log p_Z(G^{-1}(x))+\log\left|\det\frac{\partial G^{-1}}{\partial x}\right|.
\)

Composition: $z_K=G(z_0)=f_K\circ\cdots\circ f_1(z_0)$, 
$\log p_X(z_K)=\log p_Z(z_0)-\sum_k \log |\det \frac{\partial f_k}{\partial z_k}|$. 
Flow design needs invertible layer, efficient inverse, efficient compute determinant of Jacobi. 

\key{Coupling}: split $x=(x_a,x_b)$, to $(y_a,y_b)$, 
\(
y_a=x_a,\quad y_b=x_b\odot e^{s(x_a)}+t(x_a),\quad
\log|\det J|=\sum s(x_a). 
\)
$s,t$: learnable. sometimes needs premutaion among coordinates or linear transformation.  

\key{Glow}: invertible $1\times1$ convolution mixes channels. Since the same channel-mixing matrix $W$ is applied at each spatial location,
\[
\log\left|\det\frac{d\,{\rm conv2D}(h;W)}{dh}\right|=H W_{\rm spatial}\log|\det W|.
\]
LU parameterization $W=PL(U+\operatorname{diag}(s))$ makes determinant cheap: $\log|\det W|=\sum_i\log|s_i|$. \key{Dequantization}: add uniform noise to discrete pixels so continuous density is well-defined.
\key{Continuous NF / Neural ODE}: $\frac{dz(t)}{dt}=f_\theta(z(t),t)$,
\(
\log p(z(t_1))=\log p(z(t_0))-\int_{t_0}^{t_1}\tr\left(\frac{\partial f_\theta}{\partial z(t)}\right)dt.
\)

\key{Prove a Normalizing Flow}: 1. invertible; 2. easily-tractable Jacobi matrix, 

\subsection{GMM and EM}
\key{GMM responsibility}: for $K$ mixture components,
\[
\gamma_{ij}:=p(z_i=j|x_i)=
\frac{\pi_j\N(x_i|\mu_j,\Sigma_j)}
{\sum_{\ell=1}^{K}\pi_\ell\N(x_i|\mu_\ell,\Sigma_\ell)}.
\]
\key{GMM EM lower bound}:
\[
\sum_i\log\sum_j\gamma_{ij}\frac{\pi_j\N(x_i|\mu_j,\Sigma_j)}{\gamma_{ij}}
\ge
\sum_i\sum_j\gamma_{ij}\log\frac{\pi_j\N(x_i|\mu_j,\Sigma_j)}{\gamma_{ij}}.
\]
E-step fixes parameters and maximizes over $\gamma$:
\[
\gamma^{(t+1)}=\argmax_{\gamma_{ij}\ge0,\sum_j\gamma_{ij}=1}
\sum_{i,j}\gamma_{ij}\log\pi_j^{(t)}\N(x_i|\mu_j^{(t)},\Sigma_j^{(t)})
-\sum_{i,j}\gamma_{ij}\log\gamma_{ij}.
\]
M-step fixes $\gamma$ and maximizes
\[
(\pi,\mu,\Sigma)^{(t+1)}
=\argmax_{\pi,\mu,\Sigma}\sum_{i,j}\gamma_{ij}^{(t+1)}
\log \pi_j\N(x_i|\mu_j,\Sigma_j).
\]

\subsection{Evidence Lower Bound (ELBO)}
Latent-variable likelihood $\log p_\theta(x)=\log\sum_z p_\theta(x,z)$. Jensen:
\[
\log p_\theta(x)=\log\E_q\frac{p_\theta(x,z)}{q(z)}
\ge \E_q\log p_\theta(x,z)+H(q)=\mathcal L(q,\theta).
\]
Gap:
\(
\log p_\theta(x)-\mathcal L(q,\theta)=\KL(q(z)\|p_\theta(z|x)).
\) 
Equivalently
\[
\log p_\theta(x)\ge
\int q_\phi(z|x)\log p_\theta(x|z)\,dz-\KL(q_\phi(z|x)\|p(z)),
\]
with gap $\KL(q_\phi(z|x)\|p_\theta(z|x))$. EM is coordinate ascent on this lower bound: E-step sets $q(z)=p_{\theta^{old}}(z|x)$, M-step maximizes $\E_q\log p_\theta(x,z)$.

\subsection{VAE}
Encoder $q_x(z)=\N(\mu_\theta(x),\sigma_\theta(x))$, decoder $p_\psi(x|z)=\N(\psi(z),\beta I)$, prior $p(z)$. 

ELBO:
\(
\log p(x)\ge \int q_x(z)\log p_\psi(x|z)\,dz-\KL(q_x(z)\|p(z)).
\)
With Gaussian decoder,
\[
{\rm ELBO}= -\frac1{2\beta}\E_{q_x(z)}\|\psi(z)-x\|_2^2-\KL(q_x(z)\|p(z))+C.
\]
Reparameterization:
\(
z=\mu_\theta(x)+\sigma_\theta(x)\epsilon,\quad \epsilon\sim\N(0,I).
\) 
\key{Gaussian KL}: for diagonal Gaussians $p=\N(\mu_1,\sigma_1^2)$, $q=\N(\mu_2,\sigma_2^2)$,
\[
\KL(p\|q)=\log\frac{\sigma_2}{\sigma_1}+
\frac{\sigma_1^2+(\mu_1-\mu_2)^2}{2\sigma_2^2}-\frac12.
\]
For $p(z)=\N(0,I)$,
\(
\KL(q_x(z)\|p(z))=\frac12\|\mu_\theta(x)\|^2+\frac12\sigma_\theta^2(x)-\log\sigma_\theta(x)+C.
\)

\key{Terms}: reconstruction = data fidelity; KL = latent regularization. In \key{$\beta$-VAE}, large $\beta$ enforces latent variables to be \key{less} correlated with each other:
\[
\mathcal L_{\beta\text{-VAE}}
=\E_{q_\phi(z|x)}[\log p_\theta(x\mid z)]
-\beta\,\KL(q_\phi(z\mid x)\|p(z)).
\]
\key{PixelVAE/NVAE}: hierarchical latent variables improve global coherence and likelihood.

\section{L7 Diffusion and Flow Matching}
\subsection{DDPM}
\key{Forward}: gradually add Gaussian noise,
\[
q(x_t|x_{t-1}):=\N(x_t;\sqrt{1-\beta_t}\,x_{t-1},\beta_tI),\quad
\alpha_t=1-\beta_t,\quad \bar\alpha_t=\prod_{s=1}^t\alpha_s.
\]
\[
q(x_t|x_0)=\N(x_t;\sqrt{\bar\alpha_t}x_0,(1-\bar\alpha_t)I),\quad
x_t=\sqrt{\bar\alpha_t}x_0+\sqrt{1-\bar\alpha_t}\epsilon.
\]
\key{Reverse}: learn denoising transitions,
\[
p_\theta(x_{0:T})=p(x_T)\prod_{t=1}^Tp_\theta(x_{t-1}|x_t),\quad
p(x_T)=\N(0,I),
\]
\[
p_\theta(x_{t-1}|x_t)=\N(x_{t-1};\mu_\theta(x_t,t),\sigma_t^2I).
\]
\key{Loss / ELBO identity}: reuse the Gaussian-product identity
\[
\log q(x_t|x_{t-1})+\log q(x_{t-1}|x_0)
=\log q(x_{t-1}|x_t,x_0)+\log q(x_t|x_0).
\]
The generic VAE-style decomposition is
\[
-\log p(x)=-\int q_\phi(z)\log p(x|z)\,dz
+\KL(q_\phi(z)\|p(z))-\KL(q_\phi(z)\|p(z|x)).
\]
\key{True posterior}:
\[
q(x_{t-1}|x_t,x_0)=\N(x_{t-1};\tilde\mu_t(x_t,x_0),\tilde\beta_tI),
\quad
\tilde\beta_t=\frac{1-\bar\alpha_{t-1}}{1-\bar\alpha_t}\beta_t,
\]
\[
\tilde\mu_t=
\frac{\sqrt{\bar\alpha_{t-1}}\beta_t}{1-\bar\alpha_t}x_0+
\frac{\sqrt{\alpha_t}(1-\bar\alpha_{t-1})}{1-\bar\alpha_t}x_t.
\]
Gaussian KL reduces each timestep to mean matching:
\[
L_{t-1}=\E_q\left[\frac1{2\sigma_t^2}
\|\tilde\mu_t(x_t,x_0)-\mu_\theta(x_t,t)\|^2\right]+C.
\]
\key{Noise prediction}: since
\[
\tilde\mu_t=
\frac1{\sqrt{\alpha_t}}\left(x_t-\frac{\beta_t}{\sqrt{1-\bar\alpha_t}}\epsilon\right),
\quad
\mu_\theta=
\frac1{\sqrt{\alpha_t}}\left(x_t-\frac{\beta_t}{\sqrt{1-\bar\alpha_t}}\epsilon_\theta(x_t,t)\right),
\]
\[
L_{t-1}=\E_{x_0,\epsilon}\left[
\frac{\beta_t^2}{2\sigma_t^2\alpha_t(1-\bar\alpha_t)}
\|\epsilon-\epsilon_\theta(x_t,t)\|^2\right]+C.
\]
Training input is $(x_t,t)$ and target is noise $\epsilon$; $\epsilon_\theta(x_t,t)$ is usually a U-Net with time encoding. \key{Sampling} from $x_T\sim\N(0,I)$:
\[
x_{t-1}=
\frac1{\sqrt{\alpha_t}}\left(x_t-\frac{\beta_t}{\sqrt{1-\bar\alpha_t}}\epsilon_\theta(x_t,t)\right)+\sigma_tz,\quad z\sim\N(0,I).
\]

\subsection{Score-Based Models}
\lowterm{Score}: learn score $\nabla_x \log p(x)$. \lowterm{Langevin dynamics}: use only score. $x_{t+1}=x_t+\eta \nabla_x \log p(x_t)+\sqrt{2\eta}\epsilon_t, \epsilon_t\sim \N(0,I)$. continuous form: $\dd x_t=\nabla_x \log p(x_t)\dd t+\sqrt{2}\dd w_t,, \dd w\sim\N(0,d_t)$. 


\lowterm{NCSN}: noise conditional score network. 
$x_i=x_0+\sigma_i\epsilon$, learn $s_\theta(x_i,\sigma_i)$ to estimate $\nabla_x \log p(x_i)$.  

$x_t=\alpha_tx_0+\sigma_t\epsilon$, score $s_\theta(x_t,t)$? MSE loss $L=\E_{x_t}\|s_\theta(x_t,t)-\nabla_{x_t}\log p(x_t)\|^2$. $\E_{x_t}[s_\theta(x_t,t)^\top \nabla_{x_t}\log p(x_t)]= \E_{x_t}[\nabla\cdot s_\theta(x_t,t)]$, using $\nabla\cdot(sp)=(\nabla\cdot s)p+s^\top \cdot\nabla p$


\lowterm{Denoising score matching}: 
$\nabla_{x_t}\log p(x_t)=\E_{x_0\sim p(\cdot|x_t)}[\nabla_{x_t}\log p(x_t|x_0)]$
 
$\E_{x_t}\|s_\theta(x_t,t)-\nabla_{x_t}\log p(x_t|x_0)\|^2$, 

\[
L_{DSM}=\E_{x_0,t,x_t}\|s_\theta(x_t,t)-\nabla_{x_t}\log q(x_t|x_0)\|^2,
\]
where $x_t\sim q(x_t|x_0)$ is tractable. If $x_t=\alpha_t x_0+\sigma_t\epsilon$, then $\nabla_{x_t}\log q(x_t|x_0)=-\frac{x_t-\alpha_t x_0}{\sigma_t^2}=-\frac{\epsilon}{\sigma_t}$. Let $s_\theta(x_t,t)=-\frac{\epsilon_\theta(x_t,t)}{\sigma_t}$; loss matches DDPM.


SDE:
\(
dx=f(x,t)dt+g(t)dW_t,
\)

reverse: $dx=[f(x,t)-g(t)^2\nabla_x\log p_t(x)]dt+g(t)d\bar W_t$.

A equivalent ODE (same distribution when same $t$):
Probability-flow ODE:
\(
dx=\left[f(x,t)-\frac12g(t)^2\nabla_x\log p_t(x)\right]dt.
\)
Likelihood along ODE:
\(
\log p_0(x_0)=\log p_T(x_T)+\int_0^T\tr\left(\frac{\partial f_\theta}{\partial x_t}\right)dt.
\)

Inference: use learned score to solve SDE/ODE using a solver. Hence \#sample-steps in training and inference can be different. 

DDIM: sample according to schedule 


\subsection{DPM-Solver and Applications}

\lowterm{DPM solver} 1 order DPM = DDIM. 

DPM-Solver integrates diffusion ODE with high-order exponential integrators, reducing sampling steps. \lowterm{Latent diffusion}: diffuse in autoencoder latent space, cheaper than pixel space. \key{DiT/SORA}: transformer over spacetime latent patches. \key{DALL-E 2}: CLIP prior + diffusion decoder. \key{DALL-E 3}: synthetic captions improve prompt following. \key{SDEdit}: add noise to input image then denoise under prompt. \key{ControlNet}: extra condition branch injects edges/depth/pose. \key{Logit-normal sampling}: biases timestep sampling toward informative middle region.

\subsection{Flow Matching}
Continuous Normalizing Flow (CNF): $\frac{dx_t}{dt}=u_\theta(t,x_t)$, continuous change of variables $\frac{d\log p_t(x_t)}{dt}=-\operatorname{div}(u_\theta(t,x_t))$.
If $u(t,x)$ transports $p_0$ to $p_1$, Flow Matching regresses it:
$\mathcal L_{\rm FM}=\E_{t\sim U[0,1],x\sim p_t}\|u_\theta(t,x)-u(t,x)\|^2$.
\key{CFM}: conditional path, e.g. $x_t=(1-t)x_0+tx_1$, $u_t=x_1-x_0$; objective
$\mathcal L_{\rm CFM}=\E_{t,x_1\sim q,x_t\sim p_t(x|x_1)}
\|u_\theta(t,x)-u_t(x|x_1)\|^2$.
Marginal vector field is recovered in expectation. \key{Non-uniqueness}: many vector fields induce same marginals; choice affects training variance and sampling trajectory. \key{SD3}: flow target often parameterizes velocity between noise/data under continuous timestep.

\section{L8--9 Sequence Modeling}
\subsection{Sequence Data and Tokenization}
Sequence data includes language, audio, action; images/DNA/protein can be serialized. Autoregressive factorization:
\(
p(x_{1:T})=\prod_{t=1}^T p(x_t|x_{<t}).
\)

\key{Tokenization}: text $\to$ basic units $\to$ token IDs. Byte-Pair Encoding (BPE) repeatedly merges frequent adjacent symbols; balances vocabulary size and sequence length. \key{Embedding}: token ID $i\to e_i=W_E[i]$; positional info is needed because token embeddings alone are order-free.

\subsection{Autoregressive CNN and Pixel Models}
\key{WaveNet}: causal/dilated temporal convolution expands receptive field without recurrence. \key{PixelCNN}: image likelihood
\(
p(x)=\prod_i p(x_i|x_{<i}). 
\)
under raster order; masked convolution prevents access to future pixels. Gated PixelCNN uses gates to improve expressiveness.

\subsection{RNN and LSTM}
\key{RNN}:
\(
h_t=\tanh(W_{xh}x_t+W_{hh}h_{t-1}+b_h),\quad y_t=W_{hy}h_t+b_y.
\)
BPTT multiplies many Jacobians; singular values $<1$ vanish, $>1$ explode. 

\key{LSTM}:
\(
f_t=\sigma(W_{xf}x_t+W_{hf}h_{t-1}+b_f),\quad
i_t=\sigma(W_{xi}x_t+W_{hi}h_{t-1}+b_i),
\)
\(
\tilde c_t=\tanh(W_{xc}x_t+W_{hc}h_{t-1}+b_c),\quad
c_t=f_t\odot c_{t-1}+i_t\odot\tilde c_t,
\)

\(
o_t=\sigma(W_{xo}x_t+W_{ho}h_{t-1}+b_o),\quad
h_t=o_t\odot\tanh(c_t).
\)
Additive cell path gives controllable memory; forget gate decides retention.

\subsection{Attention and Transformer}
For token $t$: $q_t=W_Qx_t$, $k_j=W_Kx_j$, $v_j=W_Vx_j$,
\(
\alpha_{tj}=\frac{\exp(q_t^\top k_j/\sqrt{d_k})}{\sum_i\exp(q_t^\top k_i/\sqrt{d_k})},\quad
y_t=\sum_j\alpha_{tj}v_j.
\)
Matrix form:
\(
\operatorname{Attn}(Q,K,V)=\softmax(QK^\top/\sqrt{d_k})V.
\)

\key{Multi-head}: parallel attention subspaces, concatenate and project. \key{Causal mask}: block future tokens. \key{Transformer block}: attention + MLP + residual + normalization. Complexity $O(T^2d)$; FlashAttention reduces memory IO, Mamba/SSM explores linear-time sequence mixing.

\subsection{Objectives and Case Studies}
\key{Objectives}: autoregressive LM, masked/denoising, contrastive predictive. \key{Recipe}: pretraining, fine-tuning, low-rank tuning. \key{GPT}: decoder-only causal Transformer; next-token pretraining gives in-context learning with scale. \key{ViT}: image patches as tokens. \key{Swin}: shifted local windows for hierarchical vision. \key{CLIP}: contrastive image-text alignment,
\[
L=-\frac12\left[\frac1N\sum_i\log\frac{e^{s_{ii}/\tau}}{\sum_j e^{s_{ij}/\tau}}+
\frac1N\sum_i\log\frac{e^{s_{ii}/\tau}}{\sum_j e^{s_{ji}/\tau}}\right].
\]
\key{SigLIP}: pairwise sigmoid loss avoids global softmax. \key{Whisper}: audio Transformer trained on large-scale speech. \key{MAR}: masked autoregressive modeling for images.

\section{L10 Modern Large Models}
\subsection{Pretraining, Scaling, Prompting}
\key{Pipeline}: $\text{Pretraining}\to\text{Base GPT}\to\text{SFT}\to\text{RLHF}\to\text{ChatGPT-like}$. Pretraining builds capability by next-token prediction; post-training reshapes behavior into instruction-following. Base GPT tends to continue text; chat model tries to answer the user.
\key{Scaling laws}: held-out loss decreases smoothly as model size $N$, data $D$, and compute $C$ grow, approximately by power law. The important lesson is bottleneck balance: scale parameters, tokens, and compute together. Useful heuristics: data should grow with model size ($D\propto N^{0.74}$); compute-optimal training spends extra compute mostly on larger $N$, somewhat on batch $B$, and barely on serial steps $S$.
\key{Prompting}: zero-shot/few-shot puts task description/examples in context without parameter updates. GPT-3 made in-context learning practical; instruction tuning makes this behavior reliable by training on natural-language tasks.

\subsection{SFT and RLHF}
\key{SFT}: learn from human demonstrations $(x,y)$ by next-token imitation:
\[
\mathcal L_{\rm SFT}(\theta)
=-\E_{\mathcal D_{\rm SFT}}\sum_t\log\pi_\theta(y_t\mid x,y_{<t}).
\]
It teaches desired response style but only uses positive examples, so it cannot directly rank two acceptable answers. \key{Reward model}: pairwise preference $y_w\succ y_l$ teaches scalar reward $r_\phi(x,y)$:
\[
P_\phi(y_w\succ y_l\mid x)=\sigma(r_\phi(x,y_w)-r_\phi(x,y_l)).
\]
\[
\mathcal L_{\rm RM}=-\E\log\sigma(r_\phi(x,y_w)-r_\phi(x,y_l)).
\]
\key{RLHF}: optimize policy reward while staying near reference/SFT model:
\[
\max_\theta\E_{x\sim\mathcal D,\ y\sim\pi_\theta(\cdot\mid x)}
\left[r_\phi(x,y)-\beta\KL(\pi_\theta(\cdot\mid x)\|\pi_{\rm ref}(\cdot\mid x))\right].
\]
\key{Why KL}: prevents reward hacking and language-quality drift. \key{Pretrain vs post-train}: pretraining compresses web text distribution; post-training optimizes human utility/preference over responses.

\subsection{Hallucination, RAG, Reasoning}
\key{Hallucination}: SFT can teach a confident answer even when the base model lacks knowledge. Preference/RLHF can reward uncertainty awareness, e.g. prefer ``I don't know'' over wrong entities. Negative samples $(x,y^+,y^-)$ are valuable because they mark bad-but-plausible behavior.
\key{RAG}: combine parametric memory in $\theta$ with retrieved non-parametric memory $z$:
\[
p(y\mid x)\approx\sum_{z\in\operatorname{TopK}(x)}
p_\eta(z\mid x)\,p_\theta(y\mid x,z).
\]
RAG helps latest knowledge without retraining, but retrieval is not the same as reasoning. \key{LRM}: reasoning model spends test-time compute on internal thinking before final answer:
\[
x\rightarrow z_{1:m}\rightarrow y.
\]
\lowterm{CoT prompting} elicits intermediate reasoning through prompt examples; reasoning RL trains the behavior with verifier reward. Thinking labels are not unique, so RL can reward final correctness without prescribing one chain. Trajectory:
\[
\tau=(z_1,\ldots,z_m,y_1,\ldots,y_T),\quad
R(x,\tau)=
\begin{cases}
+1,&V(x,y)=\text{correct},\\
-1,&V(x,y)=\text{incorrect}.
\end{cases}
\]
\key{Agent loop}: reasoning can extend to environment interaction with actions:
\[
s_t=(x,\text{memory},\text{observations}_{<t}),\quad
a_t\in\{\text{search},\text{browse},\text{code},\text{write},\ldots\}.
\]

\subsection{Multimodal Models}
\key{Grounding}: multimodal models condition on text plus image/video/audio/interface inputs:
\[
p_\theta(y\mid x_{\rm text},x_{\rm image},x_{\rm video},\ldots).
\]
\key{Flamingo}: extends GPT-style few-shot prompting to interleaved text-image/video input:
\[
p_\theta(y\mid t_1,v_1,t_2,v_2,\ldots,t_k,v_k,t_{k+1}).
\]
Architecture: frozen vision encoder + frozen LM; Perceiver Resampler compresses visual features; gated cross-attention injects vision while protecting frozen LM. \key{LLaVA}: visual instruction tuning samples $(I,x_{\rm instruction},y_{\rm assistant})$; CLIP features are projected into LM embedding space:
\[
F_I=f_{\rm CLIP}(I),\quad Z_I=F_IW,\quad [Z_I;E(x_1);\ldots;E(x_T)].
\]
\key{Two-stage training}: first align visual features by training projection $W$; then freeze vision encoder and tune projection + LM for assistant behavior. \key{Unified multimodal/BAGEL}: one decoder for understanding/generation/editing/reasoning on interleaved data $t_1,I_1,t_2,I_2,\ldots,V_1,t_k$; ViT pathway supports understanding, VAE pathway supports generation/reconstruction. Emergence order: understanding $\to$ generation $\to$ editing $\to$ reasoning.

\section{Homework}
\key{Log-sum inequality}: for $a_i,b_i\ge0$ and positive sums,
\[
\sum_i a_i\log\frac{a_i}{b_i}\ge
\left(\sum_i a_i\right)\log\frac{\sum_i a_i}{\sum_i b_i}.
\]
Two-term form:
\[
(a_1+a_2)\log\frac{a_1+a_2}{b_1+b_2}
\le
a_1\log\frac{a_1}{b_1}+a_2\log\frac{a_2}{b_2}.
\]
\key{KL joint convexity}: by applying log-sum inequality pointwise and summing over $x$,
\[
\KL(\lambda p_1+(1-\lambda)p_2\|\lambda q_1+(1-\lambda)q_2)
\le
\lambda\KL(p_1\|q_1)+(1-\lambda)\KL(p_2\|q_2).
\]
\key{Gaussian density}: 1D $x\sim\N(\mu,\sigma^2)$,
$p(x)=\frac1{\sqrt{2\pi}\sigma}\exp(-\frac{(x-\mu)^2}{2\sigma^2})$.
$n$D $x\sim\N(\mu,\Sigma)$,
$p(x)=\frac{\exp[-\frac12(x-\mu)^\top\Sigma^{-1}(x-\mu)]}{(2\pi)^{n/2}|\Sigma|^{1/2}}$.
\key{Gaussian KL}: for $\N_0=\N(\mu_0,\Sigma_0)$ and $\N_1=\N(\mu_1,\Sigma_1)$ in $\R^d$,
\begin{align*}
\KL(\N_0\|\N_1)
&=\frac12\left(
\tr(\Sigma_1^{-1}\Sigma_0)
+(\mu_1-\mu_0)^\top\Sigma_1^{-1}(\mu_1-\mu_0)
\right.\\[-1pt]
&\qquad\left.
-d+\log\frac{|\Sigma_1|}{|\Sigma_0|}
\right).
\end{align*}
\key{JSD metric}: $\sqrt{\operatorname{JSD}}$ satisfies triangle inequality:
\[
\sqrt{\operatorname{JSD}(p_1\|p_2)}
+\sqrt{\operatorname{JSD}(p_2\|p_3)}
\ge
\sqrt{\operatorname{JSD}(p_1\|p_3)}.
\]
\key{$W$--TV--JSD bound}: assume $\|z-z'\|_2\le C$. Let
\[
\alpha={\rm TV}(\mu,\nu)=\frac12\|\mu-\nu\|_1,\quad
\eta(z)=\min\{\mu(z),\nu(z)\},\quad \sum_z\eta(z)=1-\alpha.
\]
Residuals: $\mu_{\rm res}=\mu-\eta$, $\nu_{\rm res}=\nu-\eta$, $\tilde\mu=\mu_{\rm res}/\alpha$, $\tilde\nu=\nu_{\rm res}/\alpha$ if $\alpha>0$. Coupling:
\[
\gamma(z,z')=\eta(z)\mathbf1_{\{z=z'\}}+\alpha\tilde\mu(z)\tilde\nu(z').
\]
Its marginals are $\mu,\nu$, so
\[
W(\mu,\nu)\le\sum_{z,z'}\gamma(z,z')\|z-z'\|_2
\le C\alpha=C\,{\rm TV}(\mu,\nu).
\]
\key{Pinsker inequality}: standard form under natural logarithm:
\[
\KL(p\|q)\ge 2\,{\rm TV}(p,q)^2,
\qquad
{\rm TV}(p,q)\le\sqrt{\frac12\KL(p\|q)}.
\]

\end{multicols*}
\end{document}
